\chapter{Detalles Técnicos}
\label{anx:tecnico}

\section{Arquitectura de agentes}

\textbf{Seshat} llama «agente» a todo componente que envuelve una invocación estructurada a un
\ac{LLM} con reintentos y control de concurrencia. Todos comparten una clase base abstracta,
\texttt{\_BaseAgent}, y desde ahí se distinguen sólo dos formas de especialización: la primera son
las familias parametrizadas del \textit{pipeline} de extracción —los agentes de identificación y
los de resolución—, y la segunda son agentes de propósito único y diverso, todos descritos a
continuación.

Por un lado, las dos \textbf{familias parametrizadas} interponen una base intermedia para generar,
por configuración, un agente concreto por tipo de concepto o par de tipos; cada familia admite
además una variante \textbf{reflectiva} opcional que añade una segunda pasada de auto-revisión
sobre la  salida del agente primario. Por otro, los \textbf{agentes de propósito único} extienden
\texttt{\_BaseAgent} directamente, y sólo difieren entre sí en su prompt y su esquema de salida:
el de \textit{grounding} verifica el anclaje de un nodo en su cita y el de agrupación consolida
los conceptos fragmentados, mientras que el extractor de palabras clave y el generador
\textit{multi-query} sirven a la recuperación de información de la base de datos vectorial.

Cada familia se expone tras un \textit{registro} que actúa como fachada: lee la configuración,
instancia los agentes concretos habilitados y los ofrece al orquestador como una colección
uniforme, de modo que este dialoga con el registro y no con cada agente por separado. Hay uno
por familia —\texttt{IdentificationRegistry} y \texttt{ResolutionRegistry}, este último apoyado
a su vez en un registro de mismo tipo y otro de tipo cruzado—. Gracias a esta abstracción, activar
o desactivar un tipo de concepto es un cambio de configuración y no del código del orquestador.
Por otro lado, los agentes de propósito único, al no parametrizarse, no pasan por ningún registro
y se construyen directamente.

\subsection{Jerarquía de clases y registro}
\label{anx:jerarquia-agentes}

La jerarquía completa de clases, colgando toda de \texttt{\_BaseAgent}, es la siguiente:

\begin{itemize}
  \item \textbf{Familia de identificación}: base \texttt{\_BaseIdentificationAgent}, que extiende
        \texttt{\_BaseAgent} con un esquema de salida genérico parametrizado por el tipo de
        concepto:
        \begin{itemize}
          \item \texttt{DecisionIdentificationAgent}
          \item \texttt{RiskIdentificationAgent}
          \item \texttt{ActionItemIdentificationAgent}
          \item \texttt{OpenQuestionIdentificationAgent}
        \end{itemize}

        Además, \texttt{ReflectiveIdentificationAgent} envuelve a cualquiera de estos agentes para
        añadir una segunda pasada de auto-revisión sobre su salida.

  \item \textbf{Familia de resolución}: base \texttt{\_BaseResolutionAgent}, con dos sub-bases
        según si el nodo fuente y el destino son del mismo tipo de concepto o de tipos distintos.
        Por un lado, \texttt{BaseSameTypeResolutionAgent} relaciona nodos de un mismo tipo y
        parametriza los tipos de relación válidos para ese \texttt{ConceptType} con un agente
        concreto por concepto:
        \begin{itemize}
          \item \texttt{DecisionResolutionAgent}
          \item \texttt{RiskResolutionAgent}
          \item \texttt{ActionItemResolutionAgent}
          \item \texttt{OpenQuestionResolutionAgent}
        \end{itemize}

        Por el otro, \texttt{BaseCrossTypeResolutionAgent} relaciona nodos de tipos distintos, con un
        agente concreto por cada tipo fuente que fija los pares \texttt{(source\_type, target\_type)}
        que ese agente gestiona:
        \begin{itemize}
          \item \texttt{DecisionCrossTypeResolutionAgent}
          \item \texttt{RiskCrossTypeResolutionAgent}
          \item \texttt{ActionItemCrossTypeResolutionAgent}
          \item \texttt{OpenQuestionCrossTypeResolutionAgent}
        \end{itemize}

        Análogamente, \texttt{ReflectiveResolutionAgent} envuelve a cualquiera de estos agentes para
        arbitrar los ambiguos en una segunda pasada.

  \item \textbf{Agentes de propósito único}: extienden \texttt{\_BaseAgent} directamente, sin base
        intermedia:
        \begin{itemize}
          \item \texttt{GroundingAgent}: verificación del anclaje de un nodo en su cita.
          \item \texttt{GroupingAgent}: agrupación de conceptos fragmentados.
          \item \texttt{KeywordAgent}: extracción de palabras clave para la búsqueda dispersa.
          \item \texttt{MultiQueryAgent}: expansión multi-consulta para la búsqueda densa.
        \end{itemize}
\end{itemize}

Sea cual sea su función, todos los agentes anteriores cumplen el mismo contrato heredado de
\texttt{\_BaseAgent}: aportan un \textit{prompt} de sistema propio, producen una salida
estructurada conforme a un modelo Pydantic dado —cuyas violaciones de tipo se rechazan y
reintentan— y encapsulan sus llamadas al proveedor con reintentos de \textit{backoff} exponencial
y \textit{jitter}. Es este contrato común, ya mencionado en el capítulo~\ref{ch:implementacion}, el
que permite tratarlos de manera uniforme pese a la diversidad de sus propósitos.

\subsection{Agentes reflectivos}
\label{anx:agentes-reflectivos}

A pesar de todas las técnicas de diseño de prompts usadas (descritas en la sección~\ref{anx:tecnicas-prompts})
y de forzar la salida estructurada los agentes, hay veces que estos agentes "shallow" no son suficientes.
Por ello, se ha implementado un modo de \textbf{auto-revisión} (\textit{self-review}) opcional sobre las
familias de agentes parametrizadas, que añade una segunda pasada de LLM sobre la salida del agente primario.

Por una parte, el \texttt{ReflectiveIdentificationAgent} implementa un ciclo \emph{extraer → validar →
filtrar}: tras la extracción, una llamada de validación comprueba el cumplimiento lógico y semántico de
cada nodo (¿satisface las reglas de extracción? ¿coincide la descripción con la cita?); los nodos que
fallan se descartan. En caso de fallo de la validación, todos los nodos se devuelven sin filtrar,
garantizando que el modo reflectivo nunca sea más perjudicial que el modo plano.

Por otra, el \texttt{ReflectiveResolutionAgent} actúa como árbitro para los casos genuinamente ambiguos.
Sólo las entradas con \texttt{alt\_rel\_type} no nulo se envían a una llamada de desempate; las entradas
claras se omiten por completo, lo que minimiza el coste adicional. Ante cualquier fallo del árbitro, se
conserva el \texttt{rel\_type} original.

Ambos agentes utilizan el patrón \textit{subclass-and-delegate}: heredan de la misma base que
los agentes planos y delegan todas las propiedades abstractas al agente interno. El orquestador
y los registros de agentes no requieren ningún cambio para soportar el modo reflectivo.

\subsection{Matriz de relaciones por par de tipos}
\label{anx:matriz-relaciones}

La tabla~\ref{tab:rel-matrix} recoge las relaciones que un agente de resolución puede asignar
para cada par dirigido (tipo fuente~$\rightarrow$~tipo destino): de las 12 combinaciones dirigidas
entre tipos distintos, sólo 9 se habilitan.

\begin{table}[H]
  \centering
  \small
  \renewcommand{\tabularxcolumn}[1]{m{#1}}
  \begin{tabularx}{\textwidth}{>{\centering\arraybackslash}m{2.7cm}!{\vrule}*{4}{>{\centering\arraybackslash}X}}
    \toprule
    \diagbox[width=\dimexpr 2.7cm+2\tabcolsep\relax, height=1.1cm, linewidth=0.2pt]{\textbf{Fuente}}{\textbf{Destino}}
      & \texttt{decision} & \texttt{risk} & \texttt{action\_item} & \texttt{open\_question} \\
    \midrule
    \texttt{decision}
      & \texttt{supersedes}, \texttt{amends}, \texttt{conflicts\_with}
      & \texttt{mitigates} & \texttt{blocks} & \texttt{resolves} \\
    \midrule[0.02em]
    \texttt{risk}
      & \texttt{blocks} & \texttt{amends} & \texttt{blocks} & \texttt{blocks} \\
    \midrule[0.02em]
    \texttt{action\_item}
      & ---
      & \texttt{mitigates}
      & \texttt{supersedes}, \texttt{amends}, \texttt{conflicts\_with}, \texttt{blocks}, \texttt{depends\_on}
      & --- \\
    \midrule[0.02em]
    \texttt{open\_question}
      & \texttt{blocks} & --- & \texttt{blocks} & \texttt{amends}, \texttt{depends\_on} \\
    \bottomrule
  \end{tabularx}
  \caption{Relaciones permitidas por par de tipos (fuente~$\rightarrow$~destino). La diagonal
           es resolución de mismo tipo; el resto, de tipo cruzado. Un guión (---) indica que
           ningún agente gestiona ese par en la versión actual.}
  \label{tab:rel-matrix}
\end{table}

\subsection{Técnicas de diseño de prompts}
\label{anx:tecnicas-prompts}

El diseño de los \textit{prompts} de los agentes combina varias técnicas complementarias, aplicadas
en mayor o menor medida según la tarea de cada uno:

\begin{itemize}
  \item \textbf{Razonamiento en cadena} (\textit{chain-of-thought}~\cite{wei2022cot}): el
        prompt descompone la tarea en pasos intermedios secuenciales y verificables —localizar
        el fragmento relevante, evaluar las compuertas lógicas, clasificar el tipo de concepto—
        antes de emitir la salida final. Hacer explícito el razonamiento reduce los errores
        al exponer los pasos intermedios a la propia lógica del modelo.

  \item \textbf{Anclaje en cita literal} (\textit{quote-first grounding}): antes de describir
        un nodo, el agente debe localizar y copiar el fragmento literal del texto fuente que lo
        respalda. Esto reduce la alucinación al obligar al modelo a anclar cada afirmación en
        evidencia concreta antes de sintetizarla; es central en la identificación y en el propio
        agente de \textit{grounding}, cuya tarea es verificar ese anclaje.

  \item \textbf{Compuertas lógicas}: el prompt enumera condiciones necesarias y suficientes
        para emitir un nodo o una relación —por ejemplo, para que un \texttt{risk} sea válido deben
        cumplirse simultáneamente cuatro condiciones: modo de fallo concreto, implicación sustantiva
        del grupo, no totalmente resuelto en el mismo intercambio, y no un requisito sin consecuencia
        declarada. Si alguna condición no se satisface, no debe emitirse. Es la técnica más presente
        en los agentes de resolución, cuyas fronteras entre tipos de relación son sutiles.

  \item \textbf{Contraejemplos etiquetados}: el prompt incluye ejemplos de casos límite que
        \emph{no} deben producirse, con explicación de por qué. Esto delimita la frontera
        semántica entre categorías próximas —por ejemplo, \texttt{risk} frente a
        \texttt{open\_question} en identificación, o \texttt{amends} frente a \texttt{supersedes}
        en resolución.

  \item \textbf{Delimitadores XML}: en los agentes que reciben la transcripción —identificación y
        \textit{grounding}—, el contenido de la reunión y los \textit{hints} de la base de
        conocimiento se encierran en bloques \texttt{<transcript>} y \texttt{<kb\_hint>}. Las
        instrucciones del sistema indican explícitamente que cualquier texto con apariencia de
        instrucción dentro de esos bloques debe tratarse como datos y no ejecutarse, lo que mitiga
        la inyección de prompts desde el contenido de la reunión.

  \item \textbf{Salida estructurada forzada}: común a todos los agentes. Los esquemas Pydantic se
        inyectan como instrucción de herramienta en la llamada al modelo; las violaciones de tipo
        activan reintentos automáticos, lo que elimina la ambigüedad de formato y hace el
        post-procesamiento determinista.
\end{itemize}



\subsection{Caché de prompts y optimización de latencia}
\label{anx:cache-prompts}

La caché de \textit{prompts} a nivel de proveedor opera sobre el \emph{prefijo} de la petición: el
tramo inicial e idéntico de una llamada a otra, que en estos agentes es el prompt de sistema. La
primera llamada que envía ese prefijo paga una prima de escritura y deja el prefijo en la caché del
proveedor durante una ventana temporal corta; cualquier llamada posterior que reutilice un prefijo
ya visto obtiene descuento de lectura sobre esos tokens. El mecanismo se activa de forma distinta
según el proveedor —en Anthropic hay que marcar el prefijo con la cabecera \texttt{cache\_control};
OpenAI cachea automáticamente los prefijos compartidos—, pero la lógica de cuándo conviene es la misma:
sólo compensa cuando ese prefijo se reutiliza en varias llamadas \emph{sucesivas}, no concurrentes,
dentro de un mismo trabajo.

Esa condición no la cumplen todos los agentes por igual, porque las llamadas al \ac{LLM} se lanzan
en paralelo bajo un límite de concurrencia; es decir, en oleadas. La primera oleada falla la caché
—todas sus llamadas escriben a la vez, sin lecturas previas—, y sólo las oleadas posteriores, o los
reintentos, la aprovechan. De ahí que la caché rinda cuando el número de llamadas supera el límite
de concurrencia:

\begin{itemize}
  \item En la \textbf{resolución}, donde cada nodo nuevo se compara con muchos candidatos, el prompt
        de sistema —y el contexto de nodos hermanos— se reutiliza a lo largo de sucesivas oleadas.
  \item En las \textbf{variantes reflectivas}, la segunda pasada de auto-revisión reemite el mismo
        prompt de sistema inmediatamente después de la primera, garantizando una lectura de caché.
  \item En el \textbf{grounding}, sólo cuando hay más nodos que verificar que el límite de
        concurrencia; en caso contrario, la única oleada no llega a beneficiarse.
\end{itemize}

La \textbf{identificación} es el caso límite: sus agentes de concepto se lanzan todos a la vez —sin
límite de concurrencia— y cada uno usa un prompt de sistema distinto, de modo que dentro de un
trabajo no hay ningún prefijo que se reutilice y amortice la escritura. Por eso no cachea su prompt
salvo en modo reflectivo, donde la pasada de validación sí lo reemite y la caché rinde. Su
transcripción, además, se excluye siempre de la caché por la misma razón de fondo: al lanzarse
todos los agentes en paralelo, cualquier lectura fallaría.

\section{Capa de observabilidad}
\label{anx:observabilidad}

La medición del consumo y la latencia descrita en la sección~\ref{sec:observabilidad} se apoya en
dos rastreadores —\texttt{UsageTracker} y \texttt{LatencyTracker}— que se instalan como
\textit{callbacks} de LangChain. El reto de implementación es que los agentes se ejecutan de forma
concurrente: la solución evita propagar el rastreador por la firma de cada método y lo publica en
una variable de contexto (\texttt{contextvars.ContextVar}). Las tareas hijas heredan el valor de la
variable al crearse, de modo que todos los agentes lanzados dentro de una etapa acumulan en el
mismo rastreador sin cambio alguno en sus firmas.

Cada etapa del \textit{pipeline} se envuelve con dos decoradores —\texttt{track\_token\_budget} y
\texttt{track\_latency\_profile}— que, al entrar, crean el rastreador y lo fijan en la variable de
contexto y, al salir, registran los totales como métricas de MLflow (tokens por tipo, y latencia
mínima, media, mediana, p95 y máxima). El decorador de presupuesto lee sus límites de la
configuración en tiempo de llamada, avisa al alcanzar el 90\,\% del cupo y aborta la etapa al
superarlo en más de un 10\,\%; este margen tolera que varios agentes concurrentes rebasen el cupo
antes de que cualquiera de ellos pueda observarlo. Los totales de cada etapa se agregan además en
un rastreador a nivel de trabajo.

Dos fuentes de consumo no pasan por las llamadas al LLM y se contabilizan mediante envoltorios
transparentes: \texttt{TrackingEmbeddings} envuelve el modelo de \textit{embeddings} y cuenta sus
tokens de entrada, y \texttt{TrackingTranscriber} envuelve al transcriptor y contabiliza los
segundos de audio procesados. Ambos empujan su medición al rastreador activo en la variable de
contexto, integrándose en el mismo presupuesto que las llamadas al LLM.

\section{Arquitectura de evaluación}

\subsection{Estructura común de los arneses y el emparejador de identificación}
\label{anx:harnesses-eval}

Los cinco arneses comparten el mismo esqueleto de ejecución, condicionado por un detalle del marco
de evaluación: \texttt{mlflow.genai.evaluate} invoca su función de predicción de forma síncrona,
mientras que los componentes de \textbf{Seshat} —orquestador y agentes— son asíncronos. Para
salvar esa incompatibilidad, cada arnés pre-computa todas las predicciones antes de entregar el
control a MLflow: un método \texttt{\_run\_all\_predictions} recorre el corpus, ejecuta el
componente real de cada caso y almacena el resultado en una caché en memoria; la función de
predicción que ve MLflow es entonces una mera consulta a esa caché. Sobre las predicciones
cacheadas, MLflow aplica el \textit{scorer} del arnés —una función determinista, sin llamadas al
LLM—, agrega las métricas por caso y el arnés vuelca el veredicto al fichero de \textit{gate}
mediante \texttt{upsert\_gate}, que actualiza sólo el bloque del arnés en ejecución y conserva el
resto.

El emparejador de identificación merece una descripción aparte, pues es la pieza no trivial de esa
cadena: debe decidir qué nodo extraído corresponde a qué nodo esperado antes de poder contar
aciertos y fallos. Resuelve un emparejamiento bipartito \textit{greedy} restringido a pares del
mismo tipo de concepto. Para cada par candidato calcula una puntuación de solapamiento en $[0,1]$
por una de dos vías: si el nodo predicho lleva cita (\texttt{quote\_anchors}), la puntuación es la
similitud difusa (\texttt{rapidfuzz}, \textit{partial ratio}) entre la cita esperada y el fragmento
de transcripción que abarca la predicha; si no la lleva, se recurre a una similitud difusa
ponderada sobre título y descripción. Los pares se ordenan de mayor a menor puntuación y se van
reclamando de forma \textit{greedy} —cada nodo esperado y cada predicho se asignan una sola vez—;
un par sólo se acepta como acierto si supera el umbral \texttt{QUOTE\_MATCH\_THRESHOLD} de 0.70. Los
nodos predichos sin pareja cuentan como falsos positivos y los esperados sin pareja, como falsos
negativos. Sobre los pares aceptados, el \textit{scorer} puntúa además la exactitud de los campos
estructurados propios de cada tipo (difusa para \texttt{assignee}, \texttt{due}, \texttt{rationale}
y \texttt{context}; exacta para el subtipo de \texttt{risk}; solapamiento de conjuntos para
\texttt{alternatives\_considered}), señales que se registran en MLflow pero no condicionan el
\textit{gate}.

Cada arnés registra en la traza de MLflow, junto a la salida que consume el \textit{scorer}, una
vista de comparación legible: los nodos —o relaciones— esperados y predichos uno al lado del otro.
El objetivo es la depurabilidad: en la interfaz de MLflow cada caso es una fila, y un caso fallido
debe poder diagnosticarse desde esa fila sin abrir el fichero YAML del corpus. Un
post-procesador de trazas recorta además los campos más verbosos de los nodos —citas en bruto,
descripciones largas— y conserva sólo tipo, título, descripción y confianza, de modo que la
comparación quepa de un vistazo.

\subsection{Caché de evaluación y meta-scorers de calibración}
\label{anx:cache-eval}

\tightparagraph{Caché de predicciones.}
La caché de predicciones que sostiene la ejecución de los arneses es también la que permite
calibrar umbrales sin repetir llamadas al LLM. Su pieza central es \texttt{read\_or\_run}: dada la
ruta de un fichero de caché y la corrutina que produciría el resultado, devuelve el contenido
cacheado si el fichero existe —cerrando la corrutina sin ejecutarla— y, en caso contrario, ejecuta
la corrutina y persiste el resultado como JSON. El nombre del fichero se deriva de tres
componentes: el identificador del caso de corpus, una huella (\textit{fingerprint}) del agente y un
\textit{hash} de la entrada. La huella del agente incorpora su \textit{prompt} y configuración, de
modo que cualquier cambio en el agente invalida automáticamente sus entradas de caché sin afectar a
las de otras variantes. Al terminar, cada arnés no borra la caché, sino que aplica un barrido de
marcado (\texttt{sweep\_stale\_entries}) que elimina únicamente las entradas en ámbito que no se
tocaron en la ejecución —predicciones obsoletas por un cambio de \textit{prompt} o de entrada—,
preservando las válidas para ejecuciones posteriores.

\tightparagraph{Meta-scorers de calibración.}
Sobre esa caché operan los dos meta-\textit{scorers}\footnote{%
  Un meta-\textit{scorer} opera un nivel por encima de un \textit{scorer}: mientras que un
  \textit{scorer} mide la calidad a un umbral \emph{fijo}, un meta-\textit{scorer} recorre
  valores discretos del rango de validez del umbral para \emph{elegir} su mejor valor.%
} de calibración. Ambos siguen un diseño en dos
fases: primero pueblan la caché con el resultado del \textit{pipeline} para cada caso —reutilizando
las predicciones que el arnés correspondiente ya hubiera generado—, y después barren el umbral
sobre esos resultados en memoria, sin más entrada/salida ni llamadas al modelo.

El meta-\textit{scorer} de identificación barre el umbral sobre la confianza heurística en el
intervalo $[0,1]$ con paso 0.01. En cada punto reproduce qué nodos se auto-aprobarían y, aplicando
el emparejador de identificación contra la verdad de referencia, calcula dos magnitudes, de forma
agregada y por tipo de concepto: la \textbf{precisión} entre los nodos aprobados y la \textbf{cobertura}
(qué fracción de los nodos esperados se llegó a auto-aprobar; es el \textit{recall}, renombrado para
subrayar que un nodo por debajo del umbral no se pierde, sólo se difiere a revisión humana). El umbral
sugerido es el que maximiza la cobertura sujeto a que la precisión alcance un objetivo configurable
(0.95 por defecto); si ninguno lo logra, se conserva el de mayor precisión, con los empates resueltos
hacia el umbral más bajo. No se emplea F1 porque contaría cada nodo diferido como un falso
negativo, justo la lectura que el renombrado de la cobertura descarta.

El meta-\textit{scorer} de \textit{retrieval} barre el umbral de similitud con paso 0.005. Para que
un único resultado cacheado sirva a cualquier corte, las búsquedas se siembran sin filtro de
puntuación y el umbral se aplica en memoria durante el barrido. El umbral sugerido es el que
maximiza el F2\footnote{%
  $F_\beta$ es la media armónica ponderada de precisión ($P$) y \textit{recall} ($R$);
  con $\beta = 2$ se obtiene $F_2 = \dfrac{5\,P\,R}{4P + R}$.%
} macro-promediado: los casos positivos aportan su F2 —que pondera el \textit{recall}
más que la precisión— y los negativos aportan especificidad (1 si no se devuelve nada por
encima del umbral, 0 en caso contrario), lo que impide que un umbral en cero, que lo aceptaría
todo, domine la selección.

\section{Esquema de la base de datos}
\label{anx:db-schema}

El estado del sistema se reparte en dos esquemas de PostgreSQL con responsabilidades separadas:
\texttt{ops}, que guarda el estado operativo (trabajos y claves de \acs{API}), y
\texttt{knowledge\_base}, que guarda el grafo de conocimiento (nodos y relaciones). Los
\textit{embeddings} viven en una tercera tabla, gestionada por LangChain, que reside en el mismo
PostgreSQL a través de la extensión pgvector.

\tightparagraph{Esquema \texttt{ops}.}
La tabla \texttt{jobs} registra el ciclo de vida de cada \textit{pipeline} —su estado, marcas de
tiempo, el \textit{blob} de entrada y el identificador de la ejecución de MLflow— e incluye una
clave de idempotencia única (\texttt{U}) que sostiene el control de duplicados. La tabla
\texttt{api\_keys} guarda las credenciales como \textit{hashes} junto con su rol. Ambas son
independientes entre sí, como muestra la figura~\ref{fig:er-ops}.

\begin{figure}[htbp]
  \centering
  \includegraphics[width=0.7\textwidth]{er_diagram_ops.drawio.png}
  \caption{Esquema \texttt{ops}. Las dos tablas no comparten claves foráneas; \texttt{U} marca la
           clave de idempotencia única de \texttt{jobs}.}
  \label{fig:er-ops}
\end{figure}

\tightparagraph{Esquema \texttt{knowledge\_base}.}
Aquí reside el grafo (figura~\ref{fig:er-kb}). \texttt{kb\_nodes} almacena cada nodo con su
contenido, su cita de anclaje y sus metadatos; \texttt{kb\_relationships} es la tabla de aristas, y
su carácter de grafo se ve en que \texttt{source\_id} y \texttt{target\_id} son ambos claves foráneas
hacia \texttt{kb\_nodes.node\_id}. Cada arista tiene una clave primaria subrogada \texttt{rel\_id} y
una restricción de unicidad sobre la terna \texttt{(source\_id, target\_id, rel\_type)} —marcada
\texttt{U1} en la figura— que impide duplicar una misma relación.

\begin{figure}[htbp]
  \centering
  \includegraphics[width=\textwidth]{er_diagram_kb.drawio.png}
  \caption{Esquema \texttt{knowledge\_base}. \texttt{PK} marca la clave primaria; \texttt{FK1} y
           \texttt{FK2}, las dos claves foráneas de \texttt{kb\_relationships} hacia
           \texttt{kb\_nodes}; y \texttt{U1}, las columnas de la restricción de unicidad de la arista.}
  \label{fig:er-kb}
\end{figure}

\tightparagraph{Almacén vectorial.}
Los \textit{embeddings} de los nodos residen en la tabla \texttt{langchain\_pg\_embedding}, que
LangChain materializa de forma diferida en la primera conexión. Sobre ella se añade una
columna generada de \textit{tsvector} con un índice \textit{GIN} que habilita la búsqueda por palabras
clave; la búsqueda semántica usa el índice vectorial de pgvector sobre la columna de \textit{embeddings}.
El vínculo con el grafo es el \texttt{node\_id}, que se conserva en los metadatos de cada
\textit{embedding} para poder recuperar el nodo correspondiente.

\section{Referencia de la API REST}
\label{anx:api-reference}

La tabla~\ref{tab:api-endpoints} agrupa los \textit{endpoints} por router: salvo \texttt{/admin} —autenticado
con la clave raíz— y \texttt{health}, todos exigen una clave de \acs{API} válida y el rol indicado
en la sección~\ref{sec:orquestacion}.

\begin{table}[H]
  \centering
  \begin{tabularx}{\textwidth}{@{}l l >{\raggedright\arraybackslash}X@{}}
  \toprule
  \textbf{Método} & \textbf{Ruta} & \textbf{Descripción} \\
  \midrule
  \multicolumn{3}{@{}l}{\textbf{\texttt{/jobs} — ciclo de vida del \textit{pipeline}}} \\
  \addlinespace
  \texttt{POST}   & \texttt{/jobs}                     & Crear un \textit{pipeline} a partir de un archivo de entrada \\
  \texttt{GET}    & \texttt{/jobs}                     & Listar \textit{pipelines} con filtros \\
  \texttt{GET}    & \texttt{/jobs/\{id\}}              & Consultar el estado de un \textit{pipeline} \\
  \texttt{GET}    & \texttt{/jobs/\{id\}/results}      & Recuperar los nodos extraídos \\
  \texttt{GET}    & \texttt{/jobs/\{id\}/transcript/excerpt} & Extraer un fragmento de la transcripción \\
  \texttt{POST}   & \texttt{/jobs/\{id\}/approve}      & Enviar decisiones de revisión y disparar la escritura \\
  \texttt{POST}   & \texttt{/jobs/\{id\}/retry}        & Reintentar un \textit{pipeline} fallido \\
  \addlinespace
  \multicolumn{3}{@{}l}{\textbf{\texttt{/graph} — consulta de la base de conocimiento}} \\
  \addlinespace
  \texttt{GET}    & \texttt{/graph}                    & Consultar nodos por filtros \\
  \texttt{GET}    & \texttt{/graph/search}             & Búsqueda semántica sobre los nodos \\
  \texttt{GET}    & \texttt{/graph/\{id\}}             & Obtener un nodo por su identificador \\
  \texttt{GET}    & \texttt{/graph/\{id\}/neighbours}  & Vecinos directos (profundidad~1) \\
  \texttt{GET}    & \texttt{/graph/\{id\}/detail}      & Nodo junto con sus vecinos directos \\
  \texttt{GET}    & \texttt{/graph/\{id\}/impact}      & Travesía multi-salto de impacto \\
  \texttt{POST}   & \texttt{/graph/nodes}              & Añadir un nodo manualmente \\
  \texttt{POST}   & \texttt{/graph/nodes/bulk}         & Añadir nodos en lote \\
  \texttt{POST}   & \texttt{/graph/nodes/resolve}      & Resolver relaciones para nodos dados \\
  \texttt{PUT}    & \texttt{/graph/nodes/\{id\}}       & Editar un nodo \\
  \texttt{PUT}    & \texttt{/graph/nodes/\{id\}/override} & Anular contenido con justificación \\
  \texttt{DELETE} & \texttt{/graph/nodes/\{id\}}       & Eliminar un nodo \\
  \texttt{DELETE} & \texttt{/graph/nodes/bulk}         & Eliminar nodos en lote \\
  \texttt{GET}    & \texttt{/graph/relationships}      & Listar relaciones \\
  \texttt{POST}   & \texttt{/graph/relationships}      & Crear una relación \\
  \texttt{DELETE} & \texttt{/graph/relationships/\{id\}} & Eliminar una relación \\
  \addlinespace
  \multicolumn{3}{@{}l}{\textbf{\texttt{/admin} — provisión de claves (clave raíz)}} \\
  \addlinespace
  \texttt{GET}    & \texttt{/admin/api-keys}           & Listar claves de \ac{API} \\
  \texttt{POST}   & \texttt{/admin/api-keys}           & Crear una clave de \ac{API} \\
  \texttt{DELETE} & \texttt{/admin/api-keys/\{id\}}    & Revocar una clave de \ac{API} \\
  \addlinespace
  \multicolumn{3}{@{}l}{\textbf{Identidad y salud}} \\
  \addlinespace
  \texttt{GET}    & \texttt{/me}                       & Resolver la clave de \ac{API} a usuario y rol \\
  \texttt{GET}    & \texttt{/health}                   & Sondeo de salud de la \ac{API} \\
  \texttt{GET}    & \texttt{/health/components}        & Estado de las dependencias externas \\
  \bottomrule
  \end{tabularx}
  \caption{%
    \textit{Endpoints} de la \acs{API} que expone \textbf{Seshat}, agrupados por router. Todos cuelgan del
    prefijo de versión \texttt{/v1}.%
  }
  \label{tab:api-endpoints}
\end{table}

\section{Configuración completa de referencia}

La referencia completa de parámetros de configuración y sus variables de entorno se mantiene junto
al código, en \href{https://github.com/josepmafe/seshat/blob/main/docs/configuration.md}{\texttt{docs/configuration.md}},
para que evolucione con él sin quedar desactualizada en este documento.

\section{\textit{Prompts} de los agentes}

Por la misma razón, los \textit{prompts} de sistema de los agentes no se reproducen en este
documento: cada uno vive junto al agente que lo usa, en su propio script, y todos los agentes están
centralizados en
\href{https://github.com/josepmafe/seshat/tree/main/src/seshat/app/agents}{\texttt{src/seshat/app/agents/}}.
